\section{Журнал выполнения}
\begin{enumerate}
    \item В ЛР1 HTML брался через \texttt{index.php?action=render}, но такой ответ не содержит \texttt{Last-Modified}, и условный запрос по нему невозможен. Робот качает обычные страницы \texttt{/wiki/<название>}: они отдают \texttt{Last-Modified} и отвечают 304 на \texttt{If-Modified-Since}.
    \item \texttt{urllib.robotparser.read()} скачивает \texttt{robots.txt} со стандартным User-Agent \texttt{Python-urllib}, на что Wikimedia отвечает 403. По правилам robotparser это означает \enquote{всё запрещено}, и робот не скачал бы ни одной страницы. Поэтому \texttt{robots.txt} скачивается самим роботом с его User-Agent, а затем передаётся в \texttt{parse()}.
    \item На полной странице есть ссылки меню сайта (\texttt{Main\_Page}, \texttt{Special:\ldots}) и блок категорий страницы внизу (\texttt{catlinks}). Через этот блок обход Wikipedia уходил бы вверх по дереву категорий, к родителям \texttt{Category:Tourism}. Решение --- брать ссылки только внутри \texttt{\#mw-content-text} (параметр \texttt{links\_in}).
    \item Ссылки на перенаправления MediaWiki (например, \texttt{World\_Tourism\_Organization}) отдают целевую статью с кодом 200 под старым URL --- получались бы дубликаты. Решение --- сохранять документ под URL из \texttt{<link rel="canonical">} (\texttt{UN\_Tourism}) и помечать повтор как \texttt{dup}.
    \item В \texttt{robots.txt} Wikipedia запрещён путь \texttt{/w/}, а ссылки на следующие страницы большой категории имеют вид \texttt{/w/index.php?title=Category:\ldots\&pagefrom=\ldots}. Робот их не обходит, поэтому из категории берутся только первые 200 подкатегорий и 200 статей.
    \item При перепроверке Wikipedia и Wikivoyage HTML неизменённой страницы отличается от запроса к запросу: в нём есть имя сервера (\texttt{wgHostname}), время ответа (\texttt{wgBackendResponseTime}) и классы A/B-экспериментов в \texttt{<body>}. Поэтому сравнение SHA-1 после ответа 200 почти ничего не отсеивает. Изменилась ли страница, на деле решает условный запрос: 200 приходит, только если \texttt{Last-Modified} страницы сдвинулся (правка, изменение шаблона, сброс кэша), иначе 304.
    \item На страницах категорий есть \enquote{красные} ссылки на несуществующие статьи (\texttt{/wiki/<название>?action=edit\&redlink=1}). Они проходят шаблон \texttt{save}, робот их скачивает и получает 404 --- в пробном прогоне 24 таких URL со статусом \texttt{error}. Отсечь их можно, запретив \texttt{?} в шаблоне.
    \item Обход категорий уходит от темы: среди категорий 2-го уровня попалась служебная \texttt{Category:Container\_categories}, и на 3-м уровне робот обошёл сотни посторонних категорий (\texttt{4th-century\_works} и т.\,п.), --- например, 16 минут подряд (407 категорий) без единой новой статьи. Для тематического корпуса такие категории нужно исключать (например, отдельным шаблоном в конфиге).
    \item SQLite --- встраиваемая база одного процесса на одной машине, а нужна распределённая. Взята MongoDB (пример из задания) в виде набора из трёх реплик в Docker. Узлы запущены в сети хоста (\texttt{network\_mode: host}): драйвер узнаёт адреса узлов из конфигурации набора, и адреса \texttt{127.0.0.1:27017--27019} одинаково доступны и самим узлам, и роботу.
    \item Транзакция MongoDB не должна длиться долго (по умолчанию не больше 60~с), а \texttt{with\_transaction} при конфликте выполняет функцию транзакции заново. Поэтому сетевые запросы (страница, \texttt{robots.txt} для новых ссылок) делаются до транзакции, а внутри неё --- только чтение и запись в базу.
    \item Данные пробного прогона перенесены из SQLite скриптом \texttt{scripts/migrate\_sqlite.py}: 2000 документов и 59\,317 записей очереди, номер записи SQLite стал \texttt{seq}, поэтому порядок обхода сохранился. Статистика по MongoDB совпала с прежней статистикой по SQLite. С учётом сжатия база занимает около 146~МиБ на узел (HTML --- 497~МиБ).
    \item \texttt{scripts/stats.py} раньше при каждом запуске перезаписывал таблицу отчёта; теперь таблица пишется только с ключом \texttt{--tex}.
\end{enumerate}

\section{План тестирования}
Тесты (\texttt{tests/run\_tests.py}) поднимают локальный HTTP-сервер (\texttt{http.server}, поддерживает \texttt{Last-Modified} и 304) с маленьким тестовым сайтом и запускают \texttt{crawler.py} как утилиту командной строки с тестовым конфигом. Проверки идут по базе и по журналу запросов сервера. Нужен запущенный набор реплик: каждый тест пишет в свою временную базу в нём, после тестов базы удаляются. Тестовый сайт: \texttt{index.html} ссылается на \texttt{a.html}, \texttt{./b.html\#top}, \texttt{/dir/../c.html}, \texttt{d.html} и внешний сайт; \texttt{a.html} --- обратно на \texttt{/index.html} (цикл) и на \texttt{b.html}; \texttt{c.html} --- на \texttt{private/secret.html} (запрещена в \texttt{robots.txt}) и несуществующую \texttt{missing.html}; у \texttt{d.html} \texttt{canonical} указывает на \texttt{b.html}.

\begin{center}
\small
\begin{tabular}{|c|p{0.27\textwidth}|p{0.6\textwidth}|}
\hline
\textnumero & Тест & Что проверяется и зачем \\
\hline
1 & Нормализация URL & 13 пар \enquote{ссылка $\to$ ожидаемый URL}: правила разрешения относительных URL из лекции, регистр, порт, якорь, кодирование пути, порядок параметров, \texttt{mailto:}/\texttt{javascript:}. Требование \enquote{нормализованный url}. \\
\hline
2 & Обход тестового сайта & Сохранены ровно 4 документа, есть все поля (источник, сырой HTML без изменений, дата в Unix time); каждая страница скачана один раз, хотя ссылки на неё записаны по-разному; дубликат по \texttt{canonical} не сохранён; \texttt{robots.txt}, внешняя и битая ссылки обработаны. Основная функция робота. \\
\hline
3 & Аварийная остановка и продолжение & Робот с задержкой 1~с убивается \texttt{kill -9} после 2-го документа, затем запускается снова. Проверяется, что в итоге есть все документы, а сохранённые до остановки не скачивались повторно. Требование \enquote{остановить в любой момент}; \texttt{kill -9} --- худший случай, Ctrl+C мягче. \\
\hline
4 & Переобкачка без изменений & Запуск с \texttt{recrawl\_after: 0}: каждый документ перепроверен условным GET, сервер ответил 304, \texttt{html} и \texttt{fetched\_at} не изменились, \texttt{checked\_at} обновился. \\
\hline
5 & Переобкачка с изменением & Меняется \texttt{b.html}; при следующем запуске она скачана заново (200), \texttt{html} и \texttt{fetched\_at} обновлены, остальные документы --- 304 без изменений. Требование \enquote{переобкачивать, только если изменились}. \\
\hline
6 & Распределённая база & Набор реплик из трёх узлов (один основной, два вторичных); документы, записанные роботом, читаются напрямую с каждого узла. \\
\hline
\end{tabular}
\end{center}

\section{Результаты}
\subsection{Автотесты}
Пройдены все 6 тестов. Журнал прогона (\texttt{results/tests.log}; строки \texttt{|} --- вывод робота):

{\scriptsize
\VerbatimInput[breaklines, breakanywhere]{../results/tests.log}
}

\subsection{Пробный прогон}
Пробный прогон выполнен версией робота с SQLite, затем данные перенесены в MongoDB (п.~11 журнала). Таблица состояния базы ниже --- по MongoDB, она совпадает с прежней.

Конфиг \texttt{config.yaml}, по 1000 документов на источник. Прогон шёл в три запуска (журналы \texttt{results/crawl-1.log}, \texttt{crawl-2.log}, \texttt{crawl-3.log}):
\begin{enumerate}
    \item 3 октября: остановлен вручную через 57~с, в базе 38 + 16 документов.
    \item Сразу после: продолжил с места остановки, снова остановлен вручную через 853~с, в базе 526 + 172 документа.
    \item 11 октября, через 8 дней (больше \texttt{recrawl\_after} = 7 дней): сначала перепроверены все 698 документов, затем докачаны 474 + 828 новых до лимита. Время --- 5530~с.
\end{enumerate}
Ни один URL не скачан как новый дважды: после остановки робот продолжал со следующего необработанного URL.

Перепроверка в 3-м запуске:
\begin{center}
\begin{tabular}{|l|r|r|r|}
\hline
Источник & 304 (не изменился) & 200, HTML обновлён & 200, HTML тот же \\
\hline
Wikivoyage & 302 & 224 & 0 \\
\hline
Wikipedia & 42 & 130 & 0 \\
\hline
\end{tabular}
\end{center}
У 341 из 354 обновлённых документов новый \texttt{Last-Modified} позже первой обкачки: страницу правили (или меняли её шаблоны) за эти 8 дней. У 13 он раньше неё --- значит, 3 октября пришёл ещё более старый \texttt{Last-Modified}, по-видимому, копия из кэша CDN. Ответов 200 с тем же HTML не было (см. п.~6 журнала). \texttt{ETag} сайты не отдают, условный запрос идёт только по \texttt{If-Modified-Since}.

Состояние базы после прогона (\texttt{scripts/stats.py}, \texttt{results/stats.tsv}):
\begin{center}
\scriptsize
\begin{tabular}{|l|r|r|r|r|r|r|r|r|}
\hline
Источник & Док. & Без сохр. & Дубл. & Ошибок & В очереди & HTML, МиБ & Ср. док., КиБ & Глубина \\
\hline
Wikivoyage & 1000 & 90 & 44 & 0 & 18072 & 302,1 & 309,4 & 3 \\
\hline
Wikipedia & 1000 & 1893 & 6 & 24 & 37188 & 195,3 & 200,0 & 4 \\
\hline
Итого & 2000 & 1983 & 50 & 24 & 55260 & 497,4 & 254,7 & 4 \\
\hline
\end{tabular}

\end{center}
\enquote{Без сохр.} --- скачанные страницы, не сохранённые как документы: категории Wikipedia и перенаправления (статья сохранена под каноническим URL); \enquote{Дубл.} --- перенаправления на уже сохранённые статьи; \enquote{Ошибок} --- 404 на \enquote{красные} ссылки (п.~7 журнала); \enquote{В очереди} --- найденные, но ещё не скачанные URL.

Скорость --- около 2,2~с на запрос к одному источнику (0,5~с задержки, остальное --- ответ сервера и разбор страницы), источники качаются параллельно. Было 8 таймаутов, все запросы прошли со второй попытки.

\subsection{Отказ узла базы}
Скрипт \texttt{scripts/failover\_demo.py} запускает робота с \texttt{config.yaml}, через 30~с останавливает основной узел базы (\texttt{docker compose stop}), через 70~с запускает его обратно, через 110~с останавливает робота. Документы в базе проверялись больше 7 дней назад, поэтому робот занимался переобкачкой.

Журнал 12 ноября (\texttt{results/failover.log}; строки \texttt{===} --- события скрипта, строки \texttt{|} --- вывод робота):

{\scriptsize
\VerbatimInput[breaklines, breakanywhere]{../results/failover.log}
}

После остановки узла 27017 оставшиеся узлы за 10~с выбрали основным 27018. В момент отказа робот ждал ответа сервера Wikipedia (таймаут), поэтому записей во время самих выборов не было; после выборов робот без перезапуска и без ошибок базы продолжил записывать обновлённые документы уже в новый основной узел. Вернувшийся узел 27017 стал вторичным и догнал остальных: в конце на каждом из трёх узлов по 2000 документов. Ошибки \texttt{error 0} в журнале --- сетевые таймауты к Wikipedia и Wikivoyage и прерванный остановкой запрос, к базе они не относятся.

\section{Выводы}
Написан поисковый робот на Python: yaml-конфиг с секциями \texttt{db}, \texttt{logic} и списком источников, распределённая база MongoDB (набор из трёх реплик) с нормализованным URL, сырым HTML, источником и датой обкачки в Unix time. При отказе основного узла база продолжает работать на оставшихся, робот переключается на новый основной узел без перезапуска. Очередь обхода хранится в базе, поэтому робот можно остановить в любой момент, и повторный запуск продолжает с того же места --- это проверено автотестом (в том числе с \texttt{kill -9}) и на реальном прогоне. Переобкачка использует условный GET: неизменённые страницы сервер не отдаёт (304), изменённые перезаписываются.

На Wikipedia и Wikivoyage робот собрал 2000 документов (около 500~МиБ HTML). Слабые места: медленный обход (около 2~с на запрос, для 15 тыс. статей на источник --- не меньше 9 часов) и уход обхода категорий от темы, для которого нужны фильтры посторонних категорий.
\pagebreak
